\section{Uniform prices for the determinant comparison}
\label{sec:ex-stationary-determinant-prices}

The absolute one-loop term counts all fine variables in a block and
therefore has a bound proportional to $L^4$. Its comparison with the
gauge-completed fine determinant has a better bound. We prove that
bound before extracting the quadratic coefficient. Throughout this
section $L$ is odd, $\mu=L^{-1}$, and all constants are chosen from
the fixed group, patch size, bridge price and penalty coefficient,
before $L$. The fine kinetic response is an additional input: on a positive
common strip of width $\sigma$, assume its value and first two
background kernels satisfy $\|K_0\|,\|K_i\|,\|K_{ij}\|\leq A_K$
in the weighted fine one-form norms, at every admissible set of
intermediate momenta. Section~\ref{sec:ex-selected-kinetic-response}
constructs that response and its bound for the selected actions.

\subsection{Differentiated Fourier matrices}
We need the value and the first two background derivatives of a matrix.
An external leg of momentum $p$ changes an input momentum $q$ to $q+p$.
Write these three kernels as $M_0(q)$, $M_i(p;q)$ and
$M_{ij}(p,r;q)$. All momenta that occur as intermediate inputs or
outputs must lie in the specified strip. A bound on these three
kernels always uses the weighted norms at their respective input and
output momenta. It does not mean an entrywise bound. The holomorphic rear reverses
the endpoints while retaining the external momenta:
\[
 (M^*)_0(q)=M_0(-q)^T,\quad
 (M^*)_i(p;q)=M_i(p;-q-p)^T,\quad
 (M^*)_{ij}(p,r;q)=M_{ij}(p,r;-q-p-r)^T.
\]
This convention is used in all products below.

Products obey the ordinary Leibniz rule with these intermediate momenta.
Thus common bounds $A,B$ for two triples give the bound $4AB$ for
their product: there is one value term, two first-derivative terms,
and four second-derivative terms. If $M_0^{-1}$ has bound $F$, its
derivative kernels follow by differentiating $M^{-1}M=1$:
\begin{align*}
 (M^{-1})_i&=-M^{-1}M_iM^{-1},\\
 (M^{-1})_{ij}&=M^{-1}M_iM^{-1}M_jM^{-1}
 +M^{-1}M_jM^{-1}M_iM^{-1}\\
 &\quad-M^{-1}M_{ij}M^{-1}.
\end{align*}
Intermediate momenta in these formulas are determined by composition.
In particular a common inverse-jet bound is
$F+F^2A+2F^3A^2+F^2A$. These formulas use only the inverse at zero
background; they impose no uniform neighborhood of invertibility in
the full background space.

Here are the local inputs to this calculus. Let $d(U)$ be the fine
covariant difference at background $U$. On each edge its difference
from $d(1)$ is the adjoint action of the edge variable minus the
identity, followed by the fixed endpoint shift. The identity
\[
 d(U)=d(1)+(d(U)-d(1))
\]
is therefore a fixed shift operator multiplied by local smooth
functions of one fine link. Positive derivatives through order two
of these functions have a factor $\mu$: every differentiated fine
link has size at most $\mu C$, and products of more such factors
are smaller since $\mu\leq1$. Passing from Euclidean norms to a
single weighted fine norm costs at most $e^4\mu^{-1}$, by
Section~\ref{sec:ex-weighted-frame-penalty}. The factor $\mu$ pays
this cost. The constant term $d(1)$ has a uniform weighted bound
directly from its multiplier $q$ and $|q|\leq h$. Consequently
the first two background derivatives of $d(U)$ are uniformly bounded
in both placements needed for $d^*d$ and $dd^*$.

The derivative of the centered frame, denoted by $X(U)$, has bound
$C L^2$ from the joint weighted space to the ordinary fine scalar
space, through two background derivatives. Indeed these are at most
third derivatives of the frame, with one internal direction of charge
one and the remaining external directions of charge zero. The frame
calculus bounds the normalized $L^2$ output; removing that
normalization multiplies by $L^2$. The logarithm and the finite chart
changes preserve this estimate. The block constraint derivative $Q(U)$
and its coarse penalty form $H(U)$ have uniformly bounded kernels.
The actual penalty Hessian $B(U)$ has the same property: its first
two background derivatives require derivatives through order four,
with two internal directions of charge one. Their total charge is
two, exactly the $L^1$ case of the weighted penalty estimate.

These estimates apply on one common smaller strip. The external
directions are the actual stationary bridge of
Section~\ref{sec:ex-selected-stationary-linear-data}. The same finite
stencils define the kernels on every torus. Thus the preceding
bounds apply to their common holomorphic Fourier matrices, rather
than to separately chosen extensions of finite-grid matrices.

\subsection{Why the regulator has a uniform trace bound}
For a map between weighted Hilbert spaces, let $\|\cdot\|_1$ denote
the nuclear norm after transforming to orthonormal coordinates.
We use
\[
 \|ABC\|_1\leq\|A\|\|B\|_1\|C\|,
 \qquad |\operatorname{tr}B|\leq\|B\|_1.
\]
The second inequality is applied to endomorphisms with the same
input and output momentum. Both statements follow from a singular
value decomposition; the first also follows by decomposing $B$ as
a sum of rank-one operators. They introduce no dimension factor.

Set
\[
 P(U)=\mu^2(d(U)^*d(U)+\mu^2I)^{-1},\qquad
 C(U)=\mu^2P(U)^9.
\]
Here and below the star means transpose with opposite Fourier
momentum in the holomorphic formulas. On the real slice it is the
usual adjoint. The scalar inverse estimate of
Section~\ref{sec:ex-weighted-stationary-inverses}, the preceding
product and inverse rules, and the two weight conversions show that
all three kernels of $P$ have a common Euclidean bound independent
of $L$. The factor $\mu^2$ cancels the two costs $O(\mu^{-1})$.
It follows that the three kernels of $C$ have operator bound
$C_0\mu^2$.

We need the stronger nuclear bound, for which the ninth power is
useful. At zero background, $P_0$ has alias eigenvalues
$d_\mu(k_l)=\mu^2/(\omega(k_l)+\mu^2)$. On the strip of width
$1/128$, after choosing a real fundamental cube, the alias distance
$n=\operatorname{dist}_\infty(l,0)$ satisfies
\[
 |d_\mu(k_l)|\leq\frac{2}{(1+n)^2}.
\]
For $n=0$ this is the scalar denominator bound. For $n>0$, the
real alias difference has size at least $2\mu n$; the same
denominator comparison gives the displayed estimate. The number
of residues with distance $n$ is at most $(2n+1)^4$. Hence
\[
 \sum_l |d_\mu(k_l)|^3
 \leq 8\sum_{n\geq0}\frac{(2n+1)^4}{(n+1)^6}
 \leq128\sum_{n\geq0}(n+1)^{-2}<\infty.
\]
The fixed color multiplicity and the bounded Fourier conjugations
only multiply this constant. Thus $(P^3)_0$ has a uniform nuclear
bound, although $P_0$ itself need not have one.

If a triple $V$ has operator bound $a$ and its value $V_0$ has
nuclear bound $b$, the triple for $V^3$ has nuclear bound
$20a^2b$. To see this, distribute at most two derivatives among
three factors. Every term has an undifferentiated factor, which
can be estimated in nuclear norm; all other factors use operator
norm. The value, first derivative and second derivative have at
most one, three and nine terms, respectively. The larger constant
20 also accommodates the grouped product estimates above.
Apply this to $V=P^3$. We obtain
\begin{equation}\label{eq:ex-regulator-nuclear}
       \|C_0\|_1,\ \|C_i\|_1,\ \|C_{ij}\|_1
                          \leq C_1\mu^2.
\end{equation}
This argument needs only the value-level spectral calculation for
$P^3$; derivatives of its eigenvectors are not used.

\subsection{Cancellation of the fine kinetic form}
Let $b$ project the joint space onto its fine one-form part, and let
$e$ project it onto its coarse scalar part. Let $s_a$ be the fixed
map selecting the one-form component in direction $a$. Denote the
fine kinetic Hessian triple by $K$. Define
\[
 X=d^*b,\qquad T=CX_{\rm fr},\qquad F=X+T,
\]
where $X_{\rm fr}$ is the centered frame derivative denoted by
$X(U)$ above. The change of notation separates it from $d^*b$.
The three matrix triples to be compared are
\begin{align*}
 M_{\rm act}&=b^*Kb+B+F^*F,\\
 M_{\rm core}&=b^*Kb+Q^*HQ+F^*F,\\
 M_{\rm rank}&=b^*\left(K+dd^*+\sum_{a=1}^4s_a^*Cs_a\right)b+e^*e.
\end{align*}
The first uses the full penalty Hessian; the second retains its
constraint pullback. Their distinction will matter in the next
comparison of low-momentum coefficients.

The weighted norm bounds above, including the bound for $K$, give
a common operator bound $A$ for all three triples and for
$d^*d+C$. Their zero-background inverses have one common bound
$F_0$ on a smaller strip of width $\eta$, by the previous two
sections. Indeed the zero-background penalty Hessian equals its
constraint pullback, and the gauge-completed fine matrix is precisely
the auxiliary matrix treated there. Take the minimum of the finitely
many widths and the maximum of their costs.

The difference of the second and third triples is exactly
\begin{equation}\label{eq:ex-core-rank-difference}
 M_{\rm core}-M_{\rm rank}
 =Q^*HQ+T^*F+X^*T-b^*\sum_a s_a^*Cs_ab-e^*e.
\end{equation}
This is an identity at all three derivative levels. In particular
the entire $K$ triple cancels, not just its value at the origin.
It follows by expanding $F=X+T$ and using $X^*X=b^*dd^*b$.

Every term on the right has a nuclear bound independent of $L$.
For $Q^*HQ$, take the nuclear bound on the fixed-dimensional
coarse space before multiplying by $Q^*$ and $Q$; the estimate
$\|H_j\|_1\leq\dim(\text{coarse one-forms})\|H_j\|$ works
for each derivative. The same argument applies to $e^*e$.
For $T=CX_{\rm fr}$, \eqref{eq:ex-regulator-nuclear} cancels
the $L^2$ frame cost because $\mu^2L^2=1$. Multiplication by
$F$ or $X$ preserves its nuclear bound. Finally, converting the
regulator term to weighted fine norms costs at most
$e^8\mu^{-2}$, paid by its factor $\mu^2$; the four component
maps $s_a$ have fixed bounds. Leibniz's rule proves the same
estimates through order two. Thus there is $k$, independent of
$L$ and of the fine kinetic price, with
\begin{equation}\label{eq:ex-core-rank-nuclear}
 \|(M_{\rm core}-M_{\rm rank})_j\|_1\leq k,
                         \qquad j=0,1,2.
\end{equation}
The inverse bound $F_0$ and absolute bound $A$ may still depend
on the stated kinetic hypotheses. That dependence has not been
removed by this cancellation.

\subsection{The second logarithmic derivative}
For any of the invertible triples define its two-leg trace density by
\[
 \mathcal T_M(i,j;p,q)=\operatorname{tr}\left[
 M_0(q)^{-1}M_{ij}(-p,p;q)
 -M_0(q)^{-1}M_i(-p;q+p)M_0(q+p)^{-1}M_j(p;q)\right].
\]
This is the second derivative of $\log\det M$, obtained by
differentiating $D\log\det M=\operatorname{tr}(M^{-1}DM)$.
It requires no choice of a global logarithm on the strip. The
formula itself is holomorphic wherever the constant inverses exist.
Restrict both independent momenta to width $\eta/2$, so every
intermediate momentum lies in width $\eta$.

If the fiber dimension is $d$, the operator bounds give
\begin{equation}\label{eq:ex-absolute-trace-density}
 |\mathcal T_M|\leq d(F_0A+F_0^2A^2).
\end{equation}
Here $d\leq L^4(d_0+d_1)$, with $d_0,d_1$ the fixed scalar
and one-form color dimensions. This explains the absolute $L^4$
cost.

For $M=M_{\rm core}$ and $N=M_{\rm rank}$ put
\[
 D=F_0kF_0,\qquad E=DA+F_0k.
\]
The resolvent identity gives
$\|M_0^{-1}-N_0^{-1}\|_1\leq D$.
The elementary telescoping identity
$UV-U'V'=(U-U')V+U'(V-V')$ then bounds the difference of
$M_0^{-1}M_{ij}$ and $N_0^{-1}N_{ij}$ by $E$ in nuclear norm.
For the four-factor term, first compare the initial two factors,
at cost $E$; compare the next inverse, at additional cost
$F_0AD$; and finally compare the last derivative. Keeping the
remaining factors in operator norm gives
\[
 \|M_0^{-1}M_iM_0^{-1}M_j
       -N_0^{-1}N_iN_0^{-1}N_j\|_1
       \leq(EF_0+F_0AD)A+F_0AF_0k.
\]
All these products respect the displayed intermediate momenta.
Taking traces proves
\begin{equation}\label{eq:ex-relative-trace-density}
 |\mathcal T_{M_{\rm core}}-\mathcal T_{M_{\rm rank}}|
 \leq E+(EF_0+F_0AD)A+F_0AF_0k.
\end{equation}
There is no dimension factor in this estimate.

The one-loop densities are
$\frac12\mathcal T_{M_{\rm act}}-\mathcal T_{d^*d+C}$,
$\frac12\mathcal T_{M_{\rm core}}-\mathcal T_{d^*d+C}$ and
$\frac12\mathcal T_{M_{\rm rank}}-\mathcal T_{d^*d+C}$.
Their individual bounds may all be enlarged to twice the right-hand
side of \eqref{eq:ex-absolute-trace-density}. In the difference
of the latter two densities the scalar determinant cancels exactly,
leaving half the difference in
\eqref{eq:ex-relative-trace-density}. This proves the required
uniform comparison price. Identification with the thermal stationary
coefficient, and the disappearance of the remaining action-versus-core
quadratic difference, are separate steps below.
